\chapter*{Аннотация}
\addcontentsline{toc}{chapter}{Аннотация}

Получили значения вязкости воздуха: для трубки с диаметром 3.95 мм $\eta_1 = 2.0 \pm 0.2$, а для трубки с диаметром 5.05 мм $\eta_2 = 1.7 \pm 0.4$, используя формулу Пуазейля. Значения в пределах погрешностей совпали с табличным значением $\eta_{\text{табл}} =  1,832 \cdot 10^{-5} \, \text{Па}. \cdot \text{с}$. Это говорит о том, что использование формулы Пуазейля позволяет с хорошей точностью определить вязкость воздуха в тонкой трубке.

Определили коэффициенты $\beta_\text{лам} = 1.7$ для ламинарного течения и $\beta_\text{турб} = 3.0$ для турбулентного течения зависимости $Q \propto R^\beta$. Полученные коэффициенты не сошлись с теоретическими значениями $\beta_{\text{лам}_\text{теор}} = 4$,  $\beta_{\text{турб}_\text{теор}} = 2.5$. Из этого можно сделать вывод, что теоретическая модель описания турбулентного течения, с допущением, что флуктуации скорости в развитом турбулентном течении по порядку величины совпадают со средней скоростью потока $\Delta u \sim  \bar{u}$, и метод размерностей для ламинарного течения оказались неприменимы для определения зависимости расхода от радиуса трубки в условиях данного эксперимента.